% 第七章：保角映射初步（详细提纲）

\chapter{保角映射初步（选学）}

\begin{wulibj}[本章定位]
前面几章里，我们已经从不同角度认识了解析函数：它有严格的可导性条件，有漂亮的积分理论，有整齐的级数展开，还有把局部信息浓缩到留数里的能力。学到这里，学生往往会形成一个很自然的印象：复变函数似乎主要是在``算积分''、``做展开''、``判奇点''。

但复分析真正迷人的地方，还不止这些。解析函数除了能参与计算，还能在几何上真正``改变区域''。更准确地说，当 $f'(z)\neq 0$ 时，解析函数在局部上就像一次``旋转加伸缩''；而把这种局部几何性质放到具体映射构造中，就会得到保角映射这一整套思想。

对数学物理方法来说，本章最重要的意义不在于多学几个变换公式，而在于帮助学生建立这样一种认识：
\begin{enumerate}
    \item 许多二维问题真正难的地方，不是方程本身，而是边界太复杂；
    \item 保角映射的价值，在于把复杂区域化成标准区域；
    \item 一旦区域被化简，拉普拉斯方程、势函数、流函数等对象就更容易处理。
\end{enumerate}

因此，本章的主线应当写成
\[
    \begin{aligned}
        &\text{解析函数的局部几何意义}
        \Longrightarrow
        \text{典型保角映射的区域作用} \\
        &\Longrightarrow
        \text{标准区域之间的互化}
        \Longrightarrow
        \text{复杂边界的简化}
        \Longrightarrow
        \text{二维势场中的简单应用}.
    \end{aligned}
\]

这也决定了本章的写法：它不宜写成纯技巧汇编，更不宜堆成一串互不相干的公式。真正要让学生学会的，是``看区域变化''、``看边界对应''、``按目标区域设计映射''的思路。

从课程编排上看，本章仍然应标为选学。一方面，它对后续主线不是硬前置；另一方面，它和第二章的解析函数局部性质、第六章的分支选择、后面数理方程里的二维边值问题都有很强联系。也就是说，它虽然是选学章，却非常适合当作复变函数部分的几何收束。
\end{wulibj}

\begin{tishi}[写作总原则]
\begin{enumerate}
    \item 开头必须先讲``为什么要学保角映射''，最好从二维静电场、稳定温度场或平面理想流动引入，而不要直接空降定义。
    \item 每讲一个映射，都应先问清楚``它把什么区域变成什么区域''，再谈公式本身。
    \item 图示是本章的核心教学工具，尤其要多画边界、角域、半平面、圆盘、条带和映射前后的对应关系。
    \item 必须反复强调：保角映射保的是角度与局部相似性，不保长度、不保面积，也不自动保证全局一一对应。
    \item 凡是涉及 $z^\alpha$、$\sqrt{z}$、$\log z$ 的映射，都要明确提醒：必须先选定解析支，否则连映射对象都没有说清楚。
    \item 黎曼映射定理与 Schwarz--Christoffel 公式可以作为章末展望提一句，但不应写成本章主线。
    \item 例题应优先训练``如何构造映射''和``如何读取边界对应''，而不是只练代数化简。
\end{enumerate}
\end{tishi}

\section{第一节：保角映射的概念与局部几何意义}

\begin{wulibj}[本节要解决的问题]
学生在第二章已经知道：若解析函数在某点导数不为零，那么它在这一点附近有很好的几何性质。但那时多数学生只把它记成一句结论，并没有真正把``局部上像旋转加伸缩''这一层理解建立起来。本节的任务，就是把这句话讲透，并把它正式提升为``保角映射''这一章的起点。
\end{wulibj}

\subsection*{教学目标}
\begin{itemize}
    \item 让学生理解解析函数在局部上为什么近似于一次复数乘法。
    \item 让学生掌握保角映射的定义，并明确其中 $f'(z)\neq 0$ 的作用。
    \item 让学生区分``局部保角''、``全局一一映射''和``保长度''这几件不同的事。
\end{itemize}

\subsection*{建议小节安排}
\begin{enumerate}
    \item 从增量公式
    \[
        f(z_0+\Delta z)-f(z_0)=f'(z_0)\Delta z+o(\Delta z)
    \]
    重新理解局部线性化
    \item 复数乘法的几何意义：旋转与伸缩
    \item 曲线夹角在映射下为什么保持
    \item 保角映射的定义与典型非例
    \item 临界点 $f'(z_0)=0$ 时会发生什么
\end{enumerate}

\subsection*{建议使用的定义与定理}
\begin{dingliyi}{点处保角}{}
设 $w=f(z)$ 在 $z_0$ 的某邻域内解析，且 $f'(z_0)\neq 0$。若通过 $z_0$ 的任意两条光滑曲线在该点的夹角，经映射后在 $w_0=f(z_0)$ 处的夹角大小和方向都保持不变，则称 $f$ 在 $z_0$ 处保角。
\end{dingliyi}

\begin{dingli}{解析函数的局部保角性}{}
若 $f(z)$ 在 $z_0$ 处解析且 $f'(z_0)\neq 0$，则 $f$ 在 $z_0$ 处保角。更具体地说，在足够小的邻域内，它近似于把图形先旋转 $\arg f'(z_0)$，再按比例 $|f'(z_0)|$ 伸缩。
\end{dingli}

\begin{dingliyi}{区域上的保角映射}{}
若 $f$ 在区域 $D$ 内解析，处处满足 $f'(z)\neq 0$，并且把 $D$ 一一映到某区域 $G$，则称 $f$ 是从 $D$ 到 $G$ 的一个保角映射。
\end{dingliyi}

\subsection*{本节必须讲清的核心点}
\begin{itemize}
    \item 保角性的本质，是局部主导项 $f'(z_0)\Delta z$ 决定了无穷小图形的变化方式。
    \item 复数乘以 $a=\rho \mathrm{e}^{\mathrm{i}\theta}$ 的几何意义，正是伸缩率 $\rho$ 与旋转角 $\theta$；这是保角性的直观来源。
    \item 保角映射保的是角度和局部相似性，不保弧长、面积，也不保全局形状。
    \item 解析且 $f'(z_0)\neq 0$ 是局部保角的充分条件；若进一步要求全局一一对应，还需要单射性。
    \item 当 $f'(z_0)=0$ 时，局部行为会像幂函数 $w=(z-z_0)^n$ 那样把角度放大，不再是通常意义下的保角。
\end{itemize}

\subsection*{建议例题}
\begin{lizi}{$w=z^2$ 在不同点的局部行为}{}
比较函数 $w=z^2$ 在 $z_0\neq 0$ 与在 $z_0=0$ 附近的几何作用，说明为什么它在原点处不保角。
\end{lizi}

\begin{lizi}{非解析映射为什么不在本章主线内}{}
讨论 $w=\overline{z}$ 对角度的影响，说明它虽然也会把直线变成直线，却不是本章所说的保角映射。
\end{lizi}

\subsection*{建议图示}
\begin{itemize}
    \item 两条曲线在 $z$ 平面交于一点，经映射后在 $w$ 平面仍保持夹角的示意图。
    \item 把一个很小的正方形映成小菱形或小平行四边形的局部图像，用来解释``旋转加伸缩''。
    \item $w=z^2$ 在原点附近把角度加倍的示意图。
\end{itemize}

\begin{jinshi}[本节易错点]
\begin{enumerate}
    \item 不要把``保角''误解成``图像形状完全不变''；它只保无穷小层面的相似性。
    \item 不要忽略定义里 $f'(z)\neq 0$ 这个条件。
    \item 不要把局部保角自动升级成全局双射。
\end{enumerate}
\end{jinshi}

\subsection*{本节习题建议}
\begin{itemize}
    \item 基础题：判断给定解析函数在某点是否保角。
    \item 综合题：分析 $z^n$ 在不同点附近的局部角度变化。
    \item 挑战题：比较 $w=z^2$ 与 $w=\overline{z}$ 对曲线夹角和方向的影响。
\end{itemize}

\section{第二节：分式线性变换}

\begin{wulibj}[本节要解决的问题]
如果第一节解决的是``为什么解析函数会保角''，那么这一节要解决的就是``最基本、最常用的一类保角映射是什么''。答案就是分式线性变换。它既是构造标准区域映射的基础模块，也是学生第一次真正体会``圆、直线、无穷远点其实可以统一看待''的地方。
\end{wulibj}

\subsection*{教学目标}
\begin{itemize}
    \item 让学生掌握分式线性变换的定义、基本分解和几何特征。
    \item 让学生理解``广义圆''的统一观点。
    \item 让学生会使用分式线性变换把上半平面、下半平面、圆盘等标准区域互相转化。
\end{itemize}

\subsection*{建议小节安排}
\begin{enumerate}
    \item 分式线性变换的定义与可逆性
    \item 平移、旋转伸缩、反演三类基本变换
    \item 扩充复平面与广义圆的统一观点
    \item 保圆性与边界对应
    \item 三点定变换的思想及其应用
\end{enumerate}

\subsection*{建议使用的定义与定理}
\begin{dingliyi}{分式线性变换}{}
形如
\[
    w=\frac{az+b}{cz+d},
    \qquad ad-bc\neq 0
\]
的映射称为分式线性变换或莫比乌斯变换。
\end{dingliyi}

\begin{dingli}{分式线性变换的基本分解}{}
任一分式线性变换都可以分解为有限次平移、旋转伸缩和反演的复合。
\end{dingli}

\begin{dingli}{保圆性}{}
分式线性变换把广义圆映成广义圆。这里所谓广义圆，是指普通圆以及把直线视为经过无穷远点的圆。
\end{dingli}

\begin{dingli}{三点定变换}{}
在扩充复平面上，给定三个互不相同的点及其像点，就能唯一确定一个分式线性变换。
\end{dingli}

\subsection*{本节必须讲清的核心点}
\begin{itemize}
    \item 分式线性变换之所以重要，不仅因为它公式简单，更因为它几乎是所有初等保角构造中的第一块积木。
    \item 反演 $w=1/z$ 是学生最容易感到突兀的一步，必须多用图像说明它如何把``不过原点的圆''、``过原点的圆''和``直线''互相转化。
    \item 讨论圆和直线时，最好始终站在扩充复平面的视角下，这样很多结论会显得统一而自然。
    \item 分式线性变换通常既保角又全局一一，但要特别注意它在极点附近要放到扩充复平面里理解。
    \item 三点定变换不是为了背公式，而是为了训练学生通过边界点和特殊点来设计映射。
\end{itemize}

\subsection*{建议例题}
\begin{lizi}{上半平面到单位圆盘}{}
验证
\[
    w=\frac{z-\mathrm{i}}{z+\mathrm{i}}
\]
把上半平面映到单位圆盘，并说明实轴为什么对应单位圆周。
\end{lizi}

\begin{lizi}{由三点条件确定分式线性变换}{}
构造一个分式线性变换，把三个指定边界点映到 $1$、$\mathrm{i}$、$-1$，并据此判断目标区域。
\end{lizi}

\begin{lizi}{反演对圆与直线的作用}{}
分别讨论一条不过原点的直线、一条过原点的直线以及一个过原点的圆在反演下的像。
\end{lizi}

\subsection*{建议图示}
\begin{itemize}
    \item 黎曼球面与球极投影的示意图，用来解释直线为什么可以看成经过无穷远点的圆。
    \item 上半平面经 Cayley 变换映到单位圆盘的边界对应图。
    \item 反演 $w=1/z$ 作用下，圆和直线互相转化的示意图。
\end{itemize}

\begin{jinshi}[本节易错点]
\begin{enumerate}
    \item 不要把直线和圆看成完全分离的两类对象；在扩充复平面里它们应统一为广义圆。
    \item 不要忘记分式线性变换在 $z=-d/c$ 处有极点，这一点需要放到扩充复平面里理解。
    \item 不要只做代数代入而不看边界对应；本节真正重要的是区域如何变化。
\end{enumerate}
\end{jinshi}

\subsection*{本节习题建议}
\begin{itemize}
    \item 基础题：判断分式线性变换把给定直线或圆映成什么。
    \item 综合题：求把上半平面映到单位圆盘、并把指定点映到原点的变换。
    \item 挑战题：利用三点定变换思想构造满足边界条件的映射。
\end{itemize}

\section{第三节：典型初等映射及其区域作用}

\begin{wulibj}[本节要解决的问题]
学生学分式线性变换时，常会误以为保角映射的世界主要就是``分式线性变换那几招''。其实真正做题时，更常用的是若干初等函数之间的复合，例如幂函数把角域变成角域，对数把角域拉直成条带，指数再把条带卷回去。本节的任务，就是把这些最常用的初等映射整理成一条清晰主线。
\end{wulibj}

\subsection*{教学目标}
\begin{itemize}
    \item 让学生掌握平移、旋转伸缩、反演、幂函数的基本几何作用。
    \item 让学生学会通过边界变化来理解映射，而不是只看点如何代入。
    \item 为下一节的标准区域互化和复合映射构造打基础。
\end{itemize}

\subsection*{建议小节安排}
\begin{enumerate}
    \item 平移与旋转伸缩：区域整体移动和转动
    \item 反演：内外翻转与边界互换
    \item 幂函数 $w=z^n$：角度放大与角域变化
    \item 一般幂函数 $w=z^\alpha$：为何必须先选支
    \item 平移、伸缩、幂映射复合时的直观读取
\end{enumerate}

\subsection*{建议使用的定义与结论}
\begin{dingli}{幂函数对辐角的作用}{}
若 $w=z^n$，则在局部上满足
\[
    |w|=|z|^n,
    \qquad
    \arg w=n\arg z.
\]
因此它会把角域的张角按 $n$ 倍放大。
\end{dingli}

\begin{dingli}{角域到半平面的典型映射}{}
若角域为
\[
    \alpha<\arg z<\beta,
\]
则在适当选支下，幂映射
\[
    w=z^{\pi/(\beta-\alpha)}
\]
可把该角域映成一个半平面。
\end{dingli}

\subsection*{本节必须讲清的核心点}
\begin{itemize}
    \item 平移和旋转伸缩虽然简单，却几乎出现在所有构造中，不能因为``太容易''而一带而过。
    \item 反演的关键作用，是把远处拉近、把内外互换，并常常把线性边界变成圆形边界。
    \item 幂函数最重要的不是计算 $z^n$，而是看清它如何改变辐角，从而改变区域的开口。
    \item 一旦写到 $z^\alpha$、$\sqrt{z}$ 等函数，就必须点明：它们不是天然单值函数，必须借助第六章的解析支思想。
    \item 本节应引导学生开始形成``看边界线变成什么''和``看特殊点变到哪里''的习惯。
\end{itemize}

\subsection*{建议例题}
\begin{lizi}{第一象限到上半平面}{}
说明为什么
\[
    w=z^2
\]
可以把第一象限映到上半平面，并分析坐标轴分别映成什么边界。
\end{lizi}

\begin{lizi}{扇形区域到半平面}{}
把张角为 $\beta-\alpha$ 的角域映到上半平面，写出映射并解释幂指数是怎样决定的。
\end{lizi}

\begin{lizi}{反演对圆盘内外的交换}{}
讨论
\[
    w=\frac{1}{z}
\]
对单位圆内外区域的作用，并说明边界为何仍映成单位圆。
\end{lizi}

\subsection*{建议图示}
\begin{itemize}
    \item 第一象限经平方映射变成上半平面的示意图。
    \item 一般角域在幂映射下张角变化的示意图。
    \item 圆盘内外在反演下交换的示意图。
\end{itemize}

\begin{jinshi}[本节易错点]
\begin{enumerate}
    \item 不要只盯着模长变化，而忽略辐角变化；在保角映射里，角度信息往往更关键。
    \item 不要未经说明就写 $z^{1/2}$、$z^\alpha$ 的映射结果，必须先说明选的是哪一支。
    \item 不要把边界映射和区域映射混为一谈；边界像看清以后，还要再判断区域究竟落在哪一侧。
\end{enumerate}
\end{jinshi}

\subsection*{本节习题建议}
\begin{itemize}
    \item 基础题：判断基本初等映射对直线、圆和角域的作用。
    \item 综合题：把指定角域映到半平面或第一象限。
    \item 挑战题：通过若干基本映射的复合，把一个简单扇形区域映到单位圆盘。
\end{itemize}

\section{第四节：标准区域之间的常见映射}

\begin{wulibj}[本节要解决的问题]
学会若干单个映射以后，学生还缺最后一块拼图：怎样把常见标准区域真正互相打通。做题时最常见的不是``任意复杂区域''，而是半平面、单位圆盘、角域、条带、带割线区域之间的互化。本节的任务，就是把这些最常用的标准通道总结清楚。
\end{wulibj}

\subsection*{教学目标}
\begin{itemize}
    \item 让学生掌握半平面、单位圆盘、角域、条带之间的典型映射。
    \item 让学生理解对数和指数在``拉直边界''与``卷回边界''中的作用。
    \item 让学生看到复合映射时，每一步都在服务最终目标区域。
\end{itemize}

\subsection*{建议小节安排}
\begin{enumerate}
    \item 上半平面与单位圆盘的互化
    \item 角域与半平面的互化
    \item 去掉一条射线的平面与条带的关系
    \item 条带与半平面的互化
    \item 典型复合映射路线总结
\end{enumerate}

\subsection*{建议使用的定义与结论}
\begin{dingli}{去心平面或带割平面到条带的映射}{}
在适当选支下，对数映射
\[
    w=\log z
\]
会把角域映成条带，把去掉一条射线的平面映成宽度为 $2\pi$ 的水平条带。
\end{dingli}

\begin{dingli}{条带到半平面的映射}{}
指数映射
\[
    \zeta=\mathrm{e}^w
\]
可把水平条带映成角域；再与分式线性变换复合，就可进一步映到半平面或单位圆盘。
\end{dingli}

\begin{dingli}{带割平面到半平面的典型映射}{}
适当选支下，
\[
    w=\sqrt{z}
\]
可以把去掉非正实轴的平面映到右半平面，把去掉非负实轴的平面映到上半平面或下半平面。
\end{dingli}

\subsection*{本节必须讲清的核心点}
\begin{itemize}
    \item 对数的核心作用，是把``角度差''变成``虚部差''，也就是把角域拉直成条带。
    \item 指数的核心作用，则是把平行边界重新卷成射线或圆周边界。
    \item 处理带割平面时，支割选法不是附属说明，而是映射结果的一部分。
    \item 半平面与单位圆盘之所以总被当作标准区域，是因为它们边界简单、结构清楚、后续边值问题也更容易处理。
    \item 本节必须把复合映射的思想点明：很多看似复杂的映射，并不是一口气想出来的，而是几步简单映射拼出来的。
\end{itemize}

\subsection*{建议例题}
\begin{lizi}{条带到上半平面}{}
构造一个映射，把条带
\[
    0<\operatorname{Im} z<\pi
\]
映到上半平面，并说明上下边界分别映到什么位置。
\end{lizi}

\begin{lizi}{带割平面到上半平面}{}
在选定主值平方根的条件下，说明
\[
    w=\sqrt{z}
\]
如何把去掉非负实轴的平面映到上半平面。
\end{lizi}

\begin{lizi}{角域到单位圆盘的复合构造}{}
先把角域映到半平面，再把半平面映到单位圆盘，写出完整复合映射并解释每一步的作用。
\end{lizi}

\subsection*{建议图示}
\begin{itemize}
    \item 角域经对数映为水平条带的示意图。
    \item 条带经指数映为角域的逆向示意图。
    \item 半平面、圆盘、条带、角域四类标准区域之间互通关系的流程图。
\end{itemize}

\begin{jinshi}[本节易错点]
\begin{enumerate}
    \item 不要一看到对数映射就忘了说明支割。
    \item 不要只判断边界像，而不判断条带内部到底落到半平面的哪一侧。
    \item 不要把不同教材里的标准区域混写，本讲义最好统一使用上半平面和单位圆盘作为最主要的终点区域。
\end{enumerate}
\end{jinshi}

\subsection*{本节习题建议}
\begin{itemize}
    \item 基础题：写出条带与半平面、半平面与圆盘之间的标准映射。
    \item 综合题：把一个简单带割平面或角域映到单位圆盘。
    \item 挑战题：给定边界对应关系，设计一条合理的复合映射路线并说明理由。
\end{itemize}

\section{第五节：保角映射的构造思路与典型例题}

\begin{wulibj}[本节要解决的问题]
到前四节为止，学生已经见过不少映射，但还未必会真正做题。最常见的困难不是不会算，而是不知道第一步该干什么。本节的任务，就是把``如何构造映射''这件事明确总结成可操作的思路，并用完整例题带学生练会。
\end{wulibj}

\subsection*{教学目标}
\begin{itemize}
    \item 让学生掌握构造保角映射时的一般思考顺序。
    \item 让学生学会依据边界对应和特殊点归一化来选择映射。
    \item 让学生在完整例题中练习``拆步构造''而不是凭感觉猜答案。
\end{itemize}

\subsection*{建议小节安排}
\begin{enumerate}
    \item 先选目标标准区域：半平面还是单位圆盘
    \item 再看边界：是直线、圆弧、射线还是角域
    \item 用特殊点做归一化：把方便的点送到 $0$、$\infty$、$\pm 1$、$\mathrm{i}$
    \item 按基本映射逐步复合
    \item 回头验证：边界、区域、单射性、支的选择是否都正确
\end{enumerate}

\subsection*{建议使用的提示框或总结框}
\begin{tishi}[构造保角映射时的五步法]
\begin{enumerate}
    \item 先把目标区域定成最熟悉的标准区域。
    \item 观察原区域边界由哪些简单曲线组成。
    \item 找几个最关键的边界点或角点，决定它们应该映到哪里。
    \item 用平移、旋转伸缩、幂函数、对数、分式线性变换逐步拼接。
    \item 最后检查：边界是否对应正确，区域是否落在正确一侧，函数是否单值解析。
\end{enumerate}
\end{tishi}

\subsection*{本节必须讲清的核心点}
\begin{itemize}
    \item 构造映射最怕一上来就猜总公式。应鼓励学生把任务拆成若干个标准步骤。
    \item ``边界对应''是第一线索，``特殊点归一化''是第二线索；两者结合后，很多映射就不再神秘。
    \item 并不是所有构造都唯一，很多题目允许不同路线；教学重点应放在思路正确，而不是答案形式唯一。
    \item 若中间步骤涉及多值函数，一定要显式说明支的选取，否则整个构造是不完整的。
    \item 本节宜用完整解答形式，不能只写结果，因为真正难的是思考路径。
\end{itemize}

\subsection*{建议例题}
\begin{lizi}{第一象限到单位圆盘}{}
构造一个映射，把第一象限映到单位圆盘，并说明为什么可以先平方、再做 Cayley 变换。
\end{lizi}

\begin{lizi}{带割平面到单位圆盘}{}
构造一个映射，把去掉非负实轴的平面映到单位圆盘，要求写清楚支的选择与每一步的几何意义。
\end{lizi}

\begin{lizi}{指定点归一化的作用}{}
在把上半平面映到单位圆盘的所有映射中，再要求某个给定点 $z_0$ 映到圆心，写出相应构造并解释其中的自由度。
\end{lizi}

\subsection*{建议图示}
\begin{itemize}
    \item ``原区域 $\to$ 中间区域 $\to$ 标准区域'' 的三步流程图。
    \item 典型角点、边界点在复合映射中的对应链条图。
\end{itemize}

\begin{jinshi}[本节易错点]
\begin{enumerate}
    \item 不要在构造完成后忘记验证区域落在哪一侧。
    \item 不要把边界上一组点的对应关系当成对整个区域的自动证明。
    \item 不要默认复合映射一定单值；一旦出现根式或对数，就必须检查定义域与支。
\end{enumerate}
\end{jinshi}

\subsection*{本节习题建议}
\begin{itemize}
    \item 基础题：把几个简单区域映到上半平面。
    \item 综合题：把角域、条带、带割平面映到单位圆盘。
    \item 挑战题：给定若干边界点的对应条件，构造一族满足要求的映射并讨论其差别。
\end{itemize}

\section{第六节：保角映射在二维势场中的简单应用}

\begin{wulibj}[本节要解决的问题]
如果本章只停留在``区域之间可以互相映来映去''，学生很容易觉得这仍然只是几何游戏。本节的任务，就是把保角映射拉回到数学物理方法的主线中，说明它为什么能帮助处理二维势场、稳恒热传导和平面理想流动中的边值问题。
\end{wulibj}

\subsection*{教学目标}
\begin{itemize}
    \item 让学生理解保角映射为什么特别适合二维拉普拉斯方程问题。
    \item 让学生看到复杂边界化为标准边界以后，边值问题会明显简化。
    \item 让学生知道本章与后面数理方程部分的联系，但不把本节写得过重。
\end{itemize}

\subsection*{建议小节安排}
\begin{enumerate}
    \item 从二维静电场或稳恒温度场引入复势思想
    \item 调和函数与解析函数实部、虚部的联系
    \item 保角映射下边界条件的对应
    \item 楔形区域、半平面或条带上的简单边值问题
    \item 物理解释：等势线、流线与几何变换
\end{enumerate}

\subsection*{建议使用的定义与结论}
\begin{dingli}{解析函数与二维势场}{}
若复势函数
\[
    F(z)=\phi(x,y)+\mathrm{i}\psi(x,y)
\]
在区域内解析，则其势函数 $\phi$ 与流函数 $\psi$ 都满足拉普拉斯方程，并互为共轭调和函数。
\end{dingli}

\begin{dingli}{保角映射下的势函数构造}{}
若 $w=f(z)$ 是从区域 $D$ 到区域 $G$ 的保角映射，而 $F(w)$ 在 $G$ 内解析，则复合函数
\[
    F(f(z))
\]
在 $D$ 内解析。于是可借助标准区域中的已知复势，在原区域中构造对应的势函数。
\end{dingli}

\subsection*{本节必须讲清的核心点}
\begin{itemize}
    \item 保角映射之所以能进入二维势场，不是因为它神奇地``解出了方程''，而是因为解析函数与调和函数之间本来就有深刻联系。
    \item 本节重点应放在``复杂边界化简''这层思想，而不是展开大规模物理计算。
    \item 最适合示范的例子，是角域或条带区域的简单 Dirichlet 边值问题，因为映射与边界条件都很直观。
    \item 要主动提醒学生：这一方法主要服务二维问题；到了三维，保角映射就不再具有同样地位。
    \item 本节结尾最好点明：后面数理方程部分会从 PDE 角度系统处理拉普拉斯方程，而本章提供的是几何化简的另一种眼光。
\end{itemize}

\subsection*{建议例题}
\begin{lizi}{楔形区域上的稳恒温度分布}{}
把一个楔形区域映到上半平面，利用边界温度为常数的简单模型，说明如何在新变量中写出势函数，再映回原区域。
\end{lizi}

\begin{lizi}{条带区域中的线性势分布}{}
在条带区域内构造满足上下边界给定常值的调和函数，并说明为什么这一例子可以通过指数映射与角域问题联系起来。
\end{lizi}

\subsection*{建议图示}
\begin{itemize}
    \item 楔形区域映到上半平面的示意图。
    \item 原区域与像区域中的边界对应、等势线与流线的对比图。
\end{itemize}

\begin{jinshi}[本节易错点]
\begin{enumerate}
    \item 不要把保角映射误解成任何 PDE 都能用的通法，它主要服务二维拉普拉斯型问题。
    \item 不要忽略边界条件在映射下的对应关系；区域映对了，边界条件也必须一起变过去。
    \item 不要把物理解释完全丢掉，本节必须回到``数学工具如何服务物理模型''这条主线。
\end{enumerate}
\end{jinshi}

\subsection*{本节习题建议}
\begin{itemize}
    \item 基础题：写出简单复势函数的势函数与流函数。
    \item 综合题：把一个角域边值问题化到上半平面。
    \item 挑战题：比较同一个二维势场问题用分离变量与用保角映射处理时各自的优缺点。
\end{itemize}

\section{本章小结与习题配置建议}

\begin{wulibj}[小结应当完成的任务]
本章小结不能只是把``保角映射、分式线性变换、角域映射''这些标题再列一遍，而要真正帮助学生回答下面四个问题：
\begin{enumerate}
    \item 本章到底在解决什么问题；
    \item 几类典型映射之间是什么关系；
    \item 做题时如何判断该先选哪类标准区域；
    \item 本章与二维势场、后续数理方程有什么联系。
\end{enumerate}
\end{wulibj}

\subsection*{小结建议结构}
\begin{enumerate}
    \item 用一句话概括本章：``保角映射的核心价值，在于把复杂边界变成标准边界。''
    \item 用流程图串起主线：
    \[
        \begin{aligned}
            &\text{局部保角性}
            \Longrightarrow
            \text{典型初等映射}
            \Longrightarrow
            \text{标准区域互化} \\
            &\Longrightarrow
            \text{复合构造}
            \Longrightarrow
            \text{二维势场应用}.
        \end{aligned}
    \]
    \item 列出本章最该记住的 5 到 6 条结论，例如：解析且导数非零时局部保角，分式线性变换保广义圆，幂函数改变张角，对数把角域拉直成条带，Cayley 变换连通半平面与圆盘。
    \item 单独整理``做题判断思路''，提醒学生先看边界形状，再看角点与特殊点，再决定是否用幂函数、对数或分式线性变换。
    \item 用 ``jinshi'' 环境集中提醒最常见误区：忘记支的选择、只看边界不看区域、把局部保角误当全局双射、忽略边界条件随映射一起变化。
\end{enumerate}

\subsection*{习题配置建议}
\begin{itemize}
    \item 基础题：判断若干映射的保角性，求简单区域的像。
    \item 综合题：构造半平面、圆盘、角域、条带之间的映射。
    \item 提高题：含分支选择的映射构造，例如带割平面到半平面或圆盘。
    \item 应用题：把一个简单二维势场问题通过映射化为标准区域问题。
    \item 选做题：介绍黎曼映射定理的结论或给出 Schwarz--Christoffel 公式的直观背景，让学生知道本章之后还能往哪里走。
\end{itemize}

\begin{tishi}[与前后章节的联系]
这一章和前面内容至少有三条明确联系：
\begin{enumerate}
    \item 它建立在第二章关于解析函数局部几何意义的基础上；
    \item 它在涉及 $\sqrt{z}$、$z^\alpha$、$\log z$ 的映射时，需要第六章关于解析支和支割的观点；
    \item 它又会把学生自然带向后面关于二维拉普拉斯方程边值问题的思考。
\end{enumerate}

因此，本章虽然标为选学，但若写得足够自学友好，就会成为整份讲义里最能体现``数学模型、几何图像与物理问题相互打通''的一章。
\end{tishi}
